% =====================================================================
%  2bat-ud4-12.tex  —  SESSIÓ 12: Optimització amb GeoGebra
%  (sense preàmbul: s'inclou des de main.tex)
% =====================================================================
\horatitol{Sessió 12 — Optimització amb GeoGebra}%
{Treball tecnològic: visualitzar les funcions d'optimització, localitzar l'òptim a la pantalla i validar que coincideix amb el càlcul fet a mà.}

\subsection{Idea clau}

Portem unes quantes sessions resolent problemes d'optimització \textbf{a mà}: muntem la
funció, la derivem i en trobem el màxim o el mínim. Avui ho \textbf{validem amb
GeoGebra}: dibuixem la funció, localitzem l'extrem a la pantalla i comprovem que el
punt que veiem coincideix amb el que havíem calculat. El càlcul, la gràfica i la
interpretació han de \textbf{dir el mateix}.

\begin{idea}
\textbf{Què validem a GeoGebra}
\begin{itemize}[leftmargin=1.4em, itemsep=2pt]
  \item Que el \textbf{punt crític} calculat a mà ($f'(x)=0$) és, a la pantalla, el
        \textbf{punt més alt (màxim) o més baix (mínim)} de la corba.
  \item Que el \textbf{valor òptim} $f(x)$ coincideix amb l'altura de l'extrem.
\end{itemize}
\textbf{Ordres útils:}
\begin{itemize}[leftmargin=1.4em, itemsep=2pt]
  \item Escriure la funció: \texttt{f(x)=-2x\^{}2+40x}.
  \item Veure la derivada: \texttt{f'(x)} (o \texttt{Derivada(f)}).
  \item Localitzar l'extrem directament: \texttt{Extrem(f)} (o \texttt{Extremum}).
  \item Marcar un punt concret: \texttt{A=(10, f(10))}.
\end{itemize}
\textbf{La regla d'or:} la gràfica i l'anàlisi s'han de validar \textbf{mútuament}; si
no coincideixen, una de les dues té un error.
\end{idea}

\subsection{Exemples}

\begin{exemple}[validar el màxim del problema del hort]
Al problema del hort (sessió 9), l'àrea era $A(x)=40x-2x^2$ i el màxim sortia a $x=10$,
amb $A(10)=200$. Ho validem: introduïm \texttt{A(x)=40x-2x\^{}2} i \texttt{Extrem(A)}.
GeoGebra marca el punt $(10,200)$ al cim de la paràbola. \textbf{Coincideix} amb el
càlcul: el punt crític és, efectivament, el punt més alt.
\begin{center}
\begin{tikzpicture}
\begin{axis}[udaxis, width=8.4cm, height=4.4cm, xlabel={\small $x$}, ylabel={\small $A$},
    xmin=0, xmax=20, ymin=0, ymax=240, xtick={0,5,10,15,20}, ytick={100,200}]
  \addplot[udblue, domain=0:20, samples=80]{40*x-2*x^2};
  \addplot[udnavy, only marks, mark=*, mark size=1.6pt] coordinates {(10,200)};
  \node[udnavy, font=\scriptsize] at (axis cs:10,222){màxim $(10,200)$};
\end{axis}
\end{tikzpicture}

\figcap{GeoGebra confirma el màxim a $(10,200)$: el punt crític calculat a mà és el cim de la corba.}
\end{center}
\end{exemple}

\begin{exemple}[descartar la solució sense sentit, amb la pantalla]
A la caixa de cartolina (sessió 10), el volum era $V(x)=4x^3-48x^2+144x$ i
\texttt{Extrem(V)} dóna \textbf{dos} punts: $(2,128)$ i $(6,0)$. A la pantalla es veu
clarament que $x=6$ correspon a un volum \textbf{zero} (la caixa desapareix), mentre que
$x=2$ és el màxim real. La gràfica ajuda a \textbf{justificar} per què descartem $x=6$.
\end{exemple}

\subsection{Exercicis resolts}

\begin{exresolt}[validar un mínim de cost]
Has calculat que el cost mitjà $f(x)=x^2-6x+13$ té un mínim a $(3,4)$. Com ho
confirmaries a GeoGebra i què hauries de veure?
\solucio
Introdueixes \texttt{f(x)=x\^{}2-6x+13} i \texttt{Extrem(f)}. GeoGebra hauria de marcar
el punt $(3,4)$ al \textbf{fons} de la paràbola (el punt més baix). També pots escriure
\texttt{f'(x)} i comprovar que talla l'eix $X$ a $x=3$ (és a dir, $f'(3)=0$). Les tres
mirades coincideixen: càlcul, gràfica i interpretació (cost mínim de $4\eur$ amb $300$
usuaris).
\resposta{GeoGebra marca el mínim a $(3,4)$; \texttt{f'(x)} talla l'eix $X$ a $x=3$.}
\end{exresolt}

\begin{exresolt}[la pantalla detecta un error de càlcul]
A mà has dit que $f(x)=-x^2+6x$ té el màxim a $x=6$. A GeoGebra, \texttt{Extrem(f)}
marca el punt a $x=3$. Qui té raó?
\solucio
Revisem el càlcul: $f'(x)=-2x+6=0\Rightarrow x=3$, no $x=6$ (segurament s'ha confós el
terme). GeoGebra té raó: el màxim és a $(3,9)$. La gràfica ha ajudat a \textbf{detectar
i corregir} l'error.
\resposta{El màxim és a $x=3$ (no $x=6$); el càlcul a mà tenia un error.}
\end{exresolt}

\subsection{Exercicis per fer}

\noindent\textit{Full de pràctica tecnològica. Per a cada funció: anota l'òptim que has
trobat a mà, comprova'l a GeoGebra i digues si coincideix.}

\begin{exercicis}
\begin{exlist}
  \item Introdueix \texttt{I(x)=-2x\^{}2+80x} i \texttt{Extrem(I)}. Quin punt marca?
        Coincideix amb el preu òptim i l'ingrés màxim que vas calcular?
  \item Introdueix \texttt{C(x)=x\^{}2-40x+500} i localitza el mínim amb
        \texttt{Extrem(C)}. A quin $x$ és? Quin és el cost mínim?
  \item Introdueix \texttt{V(x)=4x\^{}3-48x\^{}2+144x} (la caixa) i \texttt{Extrem(V)}.
        Quins \textbf{dos} punts marca? Quin descartes i per què?
  \item Dibuixa \texttt{f(x)=x\^{}2-6x+13} i la seva derivada \texttt{f'(x)}. Comprova
        que el punt on \texttt{f'(x)} talla l'eix $X$ és justament l'$x$ del mínim de
        \texttt{f}.
  \item Completa el full de validació amb tres problemes d'optimització:
        \begin{center}
        \renewcommand{\arraystretch}{1.4}
        \begin{tabular}{p{3.4cm} p{4.4cm} p{2.4cm} p{1.6cm}}
        \toprule
        \textbf{Funció} & \textbf{Òptim a mà $(x,f(x))$} & \textbf{GeoGebra} & \textbf{Coincideix?} \\
        \midrule
         & & & \\[10pt]
         & & & \\[10pt]
         & & & \\
        \bottomrule
        \end{tabular}
        \end{center}
  \item \textbf{(Validació mútua)} En algun cas la gràfica t'ha ajudat a corregir un
        error de càlcul o a descartar una solució sense sentit? Explica què havia passat.
\end{exlist}
\end{exercicis}

% ---------------------------------------------------------------------
\fullsolucions{Sessió 12}
\begin{solucions}
\begin{sollist}
  \item \texttt{Extrem(I)} marca $(20,800)$: preu òptim $20\eur$, ingrés màxim $800\eur$.
        Coincideix amb el càlcul a mà.
  \item Mínim a $x=20$; cost mínim $C(20)=400-800+500=100$ centenars d'euros. GeoGebra el
        marca al fons de la paràbola.
  \item Marca $(2,128)$ i $(6,0)$. Es descarta $x=6$ perquè dóna volum zero (el costat de
        la base $12-2\cdot 6=0$): la caixa no existeix. L'òptim real és $x=2$.
  \item \texttt{f'(x)=2x-6} talla l'eix $X$ a $x=3$, que és exactament l'$x$ del mínim de
        \texttt{f} ($(3,4)$). Coincideix.
  \item Resposta oberta (full de validació personal). Es valora la coherència entre el
        càlcul i la visualització.
  \item Resposta oberta. Es valora descriure bé el cas i l'error o la solució descartada.
\end{sollist}
\end{solucions}
